
\setlist[description]{font=\normalfont\itshape}

%%%%%%%%%%% Định dạng cho header của bảng %%%%%%%%%%%
\renewcommand\theadfont{\bfseries}
\fvset{
	tabsize=4,					% Định tabsize 
	frame=single,				% Khung
}


\captionsetup[table]{position=top}
\captionsetup{
	format=plain,
	labelsep=period,
	labelfont=bf,
	font=small,
}

%%%%%%%%%%% Cài đặt header và footer cho toàn bộ luận văn %%%%%%%%%%%
\pagestyle{fancy}
%... then configure it.
\fancyhead{} % clear all header fields
\fancyfoot{} % clear all footer fields
\fancyfoot[C]{\thepage}
\renewcommand{\headrulewidth}{0pt}


%%%%%%%%%%% Tạo không gian cho tóm tắt nội dung %%%%%%%%%%%
\makeatletter
\newenvironment{abstractV}{
  \titlepage
  \null\vfil
  \small
  \@beginparpenalty\@lowpenalty
  \begin{center}
			{\fontsize{14}{0pt}\selectfont\textbf{TÓM TẮT NỘI DUNG} \par}
                \medskip
			{\large\bfseries \nametitle \par}% Thesis title
			{\normalsize \student \par}% Author name
			\bigskip
	\@endparpenalty\@M
  \end{center}}
{\par\vfil\null\endtitlepage}
\makeatother



%%%%%%%%%%% Tạo không gian cho tóm tắt nội dung %%%%%%%%%%%
\makeatletter
\newenvironment{abstractE}{
	\titlepage
	\null\vfil
	\small
	\@beginparpenalty\@lowpenalty
	\begin{center}
		{\fontsize{14}{0pt}\selectfont\textbf{ABSTRACT} \par}
		\medskip
		{\large\bfseries \Englishnametitle \par}% Thesis title
		{\normalsize Hung Tran-Nam \par}% Author name
		\bigskip
		\@endparpenalty\@M
\end{center}}
{\par\vfil\null\endtitlepage}
\makeatother

%%%%%%%%%%% Tạo không gian cho lời cảm ơn %%%%%%%%%%%
\makeatletter
\newenvironment{acknowledgments}{
	\titlepage
	\small
	\begin{center}
	  {\fontsize{14}{0pt}\selectfont\bfseries LỜI CẢM ƠN \vspace{30pt}\vspace{\z@}}
	\end{center}
	\quotation
  }{}
\makeatother




%%%%%%%%%%% Lệnh cài đặt môi trường định lý %%%%%%%%%%%
\theoremstyle{plain}
\theoremheaderfont{\normalfont\bfseries}
\theorembodyfont{\slshape}
\theoremseparator {.} 
\renewtheorem{theorem}{\hspace*{0.5cm}Định lý} [chapter]
\newcommand{\dly}[1]{\begin{theorem}#1\end{theorem}}

%%%%%%%%%%% Lệnh cài đặt môi trường tính chất %%%%%%%%%%%
\theoremstyle{plain}
\theoremheaderfont{\normalfont\bfseries}
\theorembodyfont{\slshape}
\theoremseparator {.} 
\renewtheorem{proposition}{\hspace*{0.5cm}Tính chất} [chapter]
\newcommand{\tch}[1]{\begin{proposition}#1\end{proposition}}

%%%%%%%%%%% Lệnh cài đặt môi trường định nghĩa %%%%%%%%%%%
\theoremstyle{plain}
\theoremheaderfont{\bfseries}
\theorembodyfont{\normalfont} 
\theoremseparator {.}         
\renewtheorem{definition}{\hspace*{0.5cm}Định nghĩa}[chapter] 
\newcommand{\dng}[1]{\begin{definition}#1\end{definition}}

%%%%%%%%%%% Lệnh cài đặt môi trường ví dụ %%%%%%%%%%%
\theoremstyle{nonumberplain} 
\theoremheaderfont{\bfseries\slshape}
\theorembodyfont{\normalfont}
\renewtheorem{example}{\hspace*{0.5cm}Ví dụ}
\newcommand{\vdu}[1]{\begin{example}#1\end{example}}

%%%%%%%%%%% Lệnh cài đặt môi trường chứng minh %%%%%%%%%%%
\theoremstyle{nonumberplain} 
\theoremheaderfont{\bfseries\slshape\small}
\theorembodyfont{\small}
\theoremsymbol{\ensuremath{_\square}}
\renewtheorem{proof}{\hspace*{0.5cm}Chứng minh}
\newcommand{\chm}[1]{\begin{proof}#1\end{proof}}

\theoremstyle{nonumberplain} 
\theoremheaderfont{\bfseries\slshape}
\theorembodyfont{\normalfont}
\renewtheorem{remark}{\hspace*{0.5cm}Nhận xét}
\newcommand{\chy}[1]{\begin{remark}#1\end{remark}}

%%%%%%%%%%% Lệnh cài đặt môi trường bổ đề %%%%%%%%%%%
\theoremstyle{plain}
\theoremheaderfont{\normalfont\bfseries}
\theorembodyfont{\slshape}
\theoremseparator {.} 
\renewtheorem{lemma}{\hspace*{0.5cm}Bổ đề} [chapter]
\newcommand{\bode}[1]{\begin{lemma}#1\end{lemma}}







%%%%%%%%%%% Định dạng các mục và tiểu mục %%%%%%%%%%%
\titleformat{\chapter}[display]
{\fontsize{14}{0pt}\selectfont\bfseries\centering}{CHƯƠNG \thechapter}{0pt}{}
\titleformat{\section}
{\fontsize{13}{0pt}\selectfont\bfseries}{\thesection}{1em}{}
\titleformat{\subsection}
{\fontsize{13}{0pt}\selectfont\bfseries\slshape}{\thesubsection}{1em}{}

\titleformat{\subsubsection}
{\fontsize{13}{0pt}\selectfont\slshape}{\thesubsection}{1em}{}

\AtBeginDocument{\renewcommand{\contentsname}{MỤC LỤC}}
\AtBeginDocument{\renewcommand{\bibname}{TÀI LIỆU THAM KHẢO}}
\AtBeginDocument{\renewcommand{\listfigurename}{DANH SÁCH HÌNH VẼ}}
\AtBeginDocument{\renewcommand{\listtablename}{DANH SÁCH BẢNG}}
\AtBeginDocument{\renewcommand{\appendixname}{PHỤ LỤC}}
\AtBeginDocument{\renewcommand{\listalgorithmname}{DANH SÁCH GIẢI THUẬT}}

